\documentclass[11pt]{article}

\usepackage[utf8]{inputenc}
\usepackage{amsmath}
\usepackage{amssymb}
\usepackage{bbm}
\usepackage{mathtools}
\usepackage{geometry}
\geometry{margin=1in}

\newcommand{\Ident}{\mathbbm{1}}
\newcommand{\ffact}[2]{#1^{\underline{#2}}}   % falling factorial x^{\underline n}

\title{Local rotations for the $L=2$ XXX SoV: closed results}
\author{}
\date{}

\begin{document}
\maketitle

\begin{abstract}
\noindent
This note records the closed, hard-verified results of the ``local rotations'' program for the
two-site ($L=2$) rational $\mathfrak{gl}(2)$ (XXX) spin chain in the companion-twist SoV framework:
the pole-cancellation theorem localizing the two-site FSoV bracket to a single site, the exact
two-rotation operatorial/functional dictionary, and its explicit assembly in terms of a
state-independent recoupling matrix. All formulas below are boxed and were verified in-notebook by
hard assertion (literal-zero residuals, not numerical tolerances) across a fixed six-representation
sweep at fresh-kernel-reproducible precision. The full derivation record, including negative results
and abandoned candidates, lives in \texttt{XXX\_CG\_L2\_V2.wb}, Sections G--K (cells 78--154); this
note reports only what closed.
\end{abstract}

\section{Setup and conventions}
\label{sec:setup}

We work with the two-site inhomogeneous XXX $\mathfrak{su}(2)$ chain in the companion-twist gauge,
exactly as set up in \texttt{Clean/XXX\_CG\_L2\_Clean.wb} (Sections A--F), which this note assumes
and does not re-derive. We restate only the conventions needed to read the formulas below without
ambiguity.

\paragraph{Sites and monodromy.} Site $\alpha=1,2$ carries the $\mathfrak{gl}(2)$ representation
$V_{\lambda_\alpha}$ (highest weight $(\lambda_\alpha,0)$, dimension $\lambda_\alpha+1$) and
inhomogeneity $\theta_\alpha$. The single-site Lax matrix is
$L^{[\lambda]}(u) = u\Ident - h\sum_{i,j} E_{ji}\otimes e_{ij}$, and the two-site monodromy is the
ordered product
\begin{equation}
  T(u) = L^{[\lambda_1]}(u-\theta_1)\, L^{[\lambda_2]}(u-\theta_2) \;\in\;
  \mathrm{End}(V_{\lambda_1}\otimes V_{\lambda_2})\otimes\mathrm{End}(\mathbb{C}^2),
  \label{eq:monodromy}
\end{equation}
site 1 acting first. This ordering is not a free choice made in this note --- it is inherited from
Section A of the Clean notebook and turns out to be structurally significant: several objects below
(the recoupling matrices $A,B$ of \S\ref{sec:recoupling}) depend on $\lambda_1$ only, a direct
signature of \eqref{eq:monodromy}'s asymmetry between the two sites.

\paragraph{Companion twist.} The twist matrix is taken in companion form,
$G=\begin{psmallmatrix}\chi_1&-\chi_2\\1&0\end{psmallmatrix}$ with eigenvalues $\kappa_1,\kappa_2$,
$\chi_1=\kappa_1+\kappa_2$, $\chi_2=\kappa_1\kappa_2$. The transfer matrix is
$t_{1,1}(u)=\mathrm{tr}_a[G\,T(u)]$, and the full fusion (Hirota) hierarchy $t_{a,s}(u)$ is built
from it in the usual way, with $t_{2,1}(u)=\chi_2\,\mathrm{qdet}\,T(u)$.

\paragraph{Baxter Q-function.} For an eigenstate labelled $(M,n)$ (weight sector $M$, degeneracy
index $n$), the Baxter Q-function used throughout is
\begin{equation}
  Q_1(u) \;=\; \kappa_1^{u/h}\, q_1(u),
  \label{eq:Q1}
\end{equation}
with $q_1$ the monic polynomial solution of the TQ equation with the twist eigenvalue $\kappa_1$
chosen as reference. \emph{$Q_1$ always carries this prefactor} --- every formula below that
involves $Q_1$ uses \eqref{eq:Q1} as written, never the bare polynomial $q_1$ alone; where this
matters (\S\ref{sec:wavefunction}) it is flagged explicitly.

\paragraph{Covectors.} $\langle x_0|_\lambda$ is the lowest-weight covector of $V_\lambda^*$ in the
Gelfand--Tsetlin basis (the SoV pseudovacuum). $\langle x^{[1]}_s|_\lambda$ (written \texttt{xL1} in
the notebook) is the length-1 (single-site) SoV covector,
\begin{equation}
  \langle x^{[1]}_s|_\lambda \;=\; \frac{\langle x_0|_\lambda\, t^{[1]}_{1,s}(\theta_1)}
  {\prod_{p=1}^{s}\bigl(\theta_1+(p-1)h-\theta_1-h\lambda\bigr)},
  \qquad s=0,\dots,\lambda,
  \label{eq:xL1}
\end{equation}
built from a \emph{standalone} single-site fused hierarchy $t^{[1]}_{1,s}(u)$ (same construction as
$t_{a,s}(u)$ above, but for one site alone). $\langle x_2(s_1,s_2)|_{\lambda_1,\lambda_2}$ (written
\texttt{x2}) is the two-site SoV covector basis, built by acting on the tensor-product pseudovacuum
$\langle x_0|_{\lambda_1}\otimes\langle x_0|_{\lambda_2}$ with the two-site fused hierarchy
$t_{1,s_1}(\theta_1)\,t_{1,s_2}(\theta_2)$ and normalizing by the Yangian-weight eigenvalue
$\nu_1[\lambda_1,\lambda_2](u)=(u-\theta_1-h\lambda_1)(u-\theta_2-h\lambda_2)$ evaluated along each
tower --- this is the corrected (\texttt{Clean}) normalization, in which each tower's shift is
carried by its own site's $\lambda$, not the other site's.

\paragraph{Bracket/measure (Section C).} For each site $i=1,2$ the FSoV measure is the rational
function
\begin{equation}
  \mu_i(u) = \frac{N_i}{\prod_{k=0}^{\lambda_1}(u-\theta_1-hk)\prod_{k=0}^{\lambda_2}(u-\theta_2-hk)}
  \Bigl(1-e^{\frac{2\pi i}{h}(u-\theta_{\bar\imath})}\Bigr), \qquad \bar\imath = 3-i,
  \label{eq:mu}
\end{equation}
with $N_i$ fixed by an explicit closed form (Section C); the bracket is the residue sum over site
$i$'s tower,
\begin{equation}
  \mathrm{br}_i[f] = 2\pi i\sum_{k=0}^{\lambda_i} \mathrm{Res}_{u=\theta_i+hk}\bigl[f(u)\mu_i(u)\bigr],
  \label{eq:br}
\end{equation}
and the two-site FSoV determinant is the $1/f(\text{node})$-normalized $2\times2$ Wronskian-type
object
\begin{equation}
  \det[\lambda_1,\lambda_2][f] \;=\; \frac{1}{f(\theta_1)\,f(\theta_2)\,(\theta_2-\theta_1)}\,
  \det\!\begin{pmatrix} \mathrm{br}_1[f] & \mathrm{br}_1[uf] \\ \mathrm{br}_2[f] & \mathrm{br}_2[uf]
  \end{pmatrix}.
  \label{eq:det}
\end{equation}
The length-1 reduction (Section F) is $\mathrm{br}_1[\lambda][f] = 2\pi i\sum_{k=0}^\lambda
\mathrm{Res}_{u=\theta_1+hk}[f(u)\mu^{[1]}(u)]$ with $\mu^{[1]}(u)=N^{[1]}/\prod_{k=0}^\lambda
(u-\theta_1-hk)$, $N^{[1]}$ fixed by forcing the leading node weight to $1$, and
\begin{equation}
  \det{}_1[\lambda][f] = \frac{\mathrm{br}_1[\lambda][f]}{f(\theta_1)}.
  \label{eq:det1}
\end{equation}
Prior established context (not re-derived here): the single-rotation functional/operatorial
dictionary
\begin{equation}
  \det{}_1[\lambda_1]\bigl[k^{u/h}Q_1(u)\bigr] \;=\;
  \Bigl(\langle x_0|_{\lambda_1}\,e^{-kE_{21}}\otimes\langle x_0|_{\lambda_2}\Bigr)\cdot\Psi_{M,n},
  \label{eq:single-dict}
\end{equation}
generator $E_{21}$, universal scaling $\varphi(k)=-k$, holds to literal symbolic zero across the
sweep below (Clean Section F).

\paragraph{Numeric verification setup.} All hard-asserts below use the notebook's fixed numeric
parameters $\theta_1=\tfrac13$, $\theta_2=\tfrac17$, $h=1$, companion twist eigenvalue
$\kappa_1=e^{i\zeta(3)}$ (30-digit precision), and the fixed six-representation sweep
$(\lambda_1,\lambda_2)\in\{(1,1),(2,1),(1,2),(2,2),(3,1),(1,3)\}$, giving $41$ eigenstates in total
($4+6+6+9+8+8$). We abbreviate $\theta_{12}:=\theta_1-\theta_2$.

\section{Pole-cancellation theorem}
\label{sec:pole-cancel}

Multiplying the test function by a polynomial with roots on the \emph{other} site's inhomogeneity
tower collapses the two-site bracket \eqref{eq:br} exactly onto the pure single-site bracket:
\begin{equation}
  \mathrm{br}_1\Bigl[P^{(2)}_{\lambda_2}(u)\,f(u)\Bigr] \;=\; C_1(\lambda_1,\lambda_2)\;\mathrm{br}_1[\lambda_1][f],
  \qquad
  P^{(2)}_{\lambda_2}(u) := \prod_{k=0}^{\lambda_2}\bigl(u-\theta_2-hk\bigr),
  \label{eq:pole-cancel}
\end{equation}
\begin{equation}
  C_1(\lambda_1,\lambda_2) \;=\; (-1)^{\lambda_2-1}\prod_{k=0}^{\lambda_2}\bigl(\theta_2+kh-\theta_1\bigr).
  \label{eq:C1}
\end{equation}
The mirror statement, cancelling site 1's tower instead, holds with $1\leftrightarrow2$:
$\mathrm{br}_2\bigl[P^{(1)}_{\lambda_1}(u) f(u)\bigr] = C_2(\lambda_1,\lambda_2)\,\mathrm{br}_1^{(2)}[\lambda_2][f]$,
$P^{(1)}_{\lambda_1}(u)=\prod_{k=0}^{\lambda_1}(u-\theta_1-hk)$,
$C_2(\lambda_1,\lambda_2)=(-1)^{\lambda_1-1}\prod_{k=0}^{\lambda_1}(\theta_1+kh-\theta_2)$, where
$\mathrm{br}_1^{(2)}[\lambda_2][\cdot]$ is the site-2 analogue of \eqref{eq:det1}'s bracket (built by
re-centering the $N^{[1]}/\mu^{[1]}/\mathrm{br}_1$ construction around $\theta_2$). Bracket
localization is thus site-symmetric: the same mechanism works at either site with the roles of
$\theta_1,\theta_2,\lambda_1,\lambda_2$ exchanged.

\emph{Verification.} \eqref{eq:pole-cancel}--\eqref{eq:C1} and their mirror were hard-asserted to
literal symbolic $0$ on a basis of monomials $u^j$ and twisted monomials $k^{u/h}u^j$
($j=0,\dots,\lambda_1+2$) across the full sweep, and numerically (residual $0$ to the working
precision) on the actual $Q_1(u)$ of every one of the 41 states, in both directions (Sections G,
cells 79--87).

\section{Global-rotation dictionary}
\label{sec:global}

Define the two-rotation operatorial overlap
\begin{equation}
  O(k_1,k_2) \;:=\; \Bigl(\langle x_0|_{\lambda_1}e^{-k_1E_{21}}\otimes\langle x_0|_{\lambda_2}e^{-k_2E_{21}}\Bigr)\cdot\Psi_{M,n}.
  \label{eq:Odef}
\end{equation}
On the diagonal $k_1=k_2=k$, the full two-site FSoV determinant \eqref{eq:det} applied to the same
twisted Q-function used in \eqref{eq:single-dict} reproduces $O$ exactly:
\begin{equation}
  \boxed{\;\det[\lambda_1,\lambda_2]\bigl[k^{u/h}Q_1(u)\bigr] \;=\; O(k,k)\;}
  \label{eq:step0}
\end{equation}
\emph{Verification:} literal symbolic zero residual over the full 41-state sweep (Section I, cells
107--108). This is the calibration anchor for everything that follows.

\section{Rotated-covector identification}
\label{sec:covector-id}

The finite unipotent rotation of the lowest-weight covector by a power of the lowering generator is
exactly proportional to the single-site SoV covector \eqref{eq:xL1}:
\begin{equation}
  \boxed{\;\langle x_0|_\lambda\, E_{21}^m \;=\; \gamma_m(\lambda)\,\langle x^{[1]}_m|_\lambda\;},
  \qquad
  \gamma_m(\lambda) \;=\; \frac{\lambda!}{(\lambda-m)!} \;=\; \ffact{\lambda}{m}
  \label{eq:gamma}
\end{equation}
(falling factorial), a pure representation-theoretic statement with no $\theta$- or site-dependence.
\emph{Verification:} exact for $\lambda=1,2,3$, all $m=0,\dots,\lambda$ (Section I, cells 109--110).

\section{Level-graded recoupling}
\label{sec:recoupling}

The naive guess that the two-site covector factorizes as a single tensor product,
$\langle x_2(s_1,m)| \propto \langle x^{[1]}_{s_1}|_{\lambda_1}\otimes\langle x^{[1]}_m|_{\lambda_2}$,
holds \emph{only} at $m=0$ (a prior, independently established Clean-notebook result). For $m\ge1$
it is replaced by a genuine \textbf{level-graded recoupling}:
\begin{equation}
  \boxed{\;\langle x_2(s_1,m)| \;=\;\!\!\sum_{\substack{s_1'+m'=s_1+m\\ 0\le s_1'\le\lambda_1,\;0\le m'\le\lambda_2}}\!\!
  A(s_1,m;s_1',m')\;\langle x^{[1]}_{s_1'}|_{\lambda_1}\otimes\langle x^{[1]}_{m'}|_{\lambda_2}\;}
  \label{eq:A-decomp}
\end{equation}
supported exactly on the total level $s_1'+m'=s_1+m$ (\emph{verified}: zero off-level coefficients
in $127$ symbolic cases, $\lambda_1,\lambda_2\le3$). Within each level block, $A$ is upper-triangular
in $s_1$ ascending order (a target $\langle x_2(s_1,m)|$ only involves basis vectors with
$s_1'\ge s_1$), and its diagonal has the closed form
\begin{equation}
  \boxed{\;A(s,m;s,m) \;=\; \frac{\ffact{(s+\theta_{12})}{m}}{\ffact{(\lambda_1+\theta_{12})}{m}}
  \;=\; \frac{(s+\theta_{12})(s+\theta_{12}-1)\cdots(s+\theta_{12}-m+1)}
             {(\lambda_1+\theta_{12})(\lambda_1+\theta_{12}-1)\cdots(\lambda_1+\theta_{12}-m+1)}\;},
  \label{eq:A-diag}
\end{equation}
worked out from small ($\lambda_1,\lambda_2\le2$) fully symbolic cases (genuinely free symbols
$\theta_A,\theta_B$, not the notebook's already-numeric $\theta_1,\theta_2$) and verified exactly
against the numeric $\theta_1,\theta_2$ over $81$ cases, $\lambda_1,\lambda_2\le3$ (Section J, cells
125--138). At $s=\lambda_1$ this gives $A=1$ for every $m$: the extremal $s_1=\lambda_1$ slice is
always unmixed.

\subsection{The inverse recoupling $B$}

Define $B:=A^{-1}$, block-by-block within each level. This is not an independent construction: it
was solved directly by linear algebra --- the mixed product-covector pairing
$\mathrm{MP}(a,b):=\bigl(\langle x^{[1]}_a|_{\lambda_1}\otimes\langle x^{[1]}_b|_{\lambda_2}\bigr)\cdot\Psi$
lies in the span of the $x_2$-pairings $W(s,t):=\langle x_2(s,t)|\cdot\Psi$ (a state-independent
statement, since \eqref{eq:A-decomp} holds at the level of covectors), so
$\mathrm{MP}(a,b)=\sum_{s+t=a+b}B(a,b;s,t)\,W(s,t)$ with $B$ state-independent; solving this
massively overdetermined linear system (every state of a representation gives one equation) against
$B=A^{-1}$ agrees to exact rational equality. $B$'s diagonal is the reciprocal of \eqref{eq:A-diag}:
\begin{equation}
  \boxed{\;B(a,b;a,b) \;=\; \frac{\ffact{(\lambda_1+\theta_{12})}{b}}{\ffact{(a+\theta_{12})}{b}}\;}
  \label{eq:B-diag}
\end{equation}
Both $A$ and $B$ were found to depend on $\lambda_1$ \emph{only} (not $\lambda_2$) for every entry
checked --- a direct consequence of the monodromy ordering \eqref{eq:monodromy}: site 1 enters the
fused hierarchy first, so the recoupling structure it generates is blind to the size of site 2's
representation, which only bounds which entries exist. The off-diagonal closed form of $A$ (hence
$B$) was not found in fully general closed form (see \S\ref{sec:closing}); it is available as exact
rationals for $\lambda_1,\lambda_2\le3$ directly in the notebook.

\section{Normalized wave-function factorization}
\label{sec:wavefunction}

The two-site covector pairs with the eigenstate exactly as a ratio of Baxter Q-function values at
the corresponding SoV lattice points:
\begin{equation}
  \boxed{\;W(s_1,s_2) \;:=\; \langle x_2(s_1,s_2)|\cdot\Psi_{M,n} \;=\;
  \frac{Q_1(\theta_1+s_1h)\,Q_1(\theta_2+s_2h)}{Q_1(\theta_1)\,Q_1(\theta_2)}\;}
  \label{eq:wavefunction}
\end{equation}
with $Q_1$ \emph{exactly as in \eqref{eq:Q1}, including its $\kappa_1^{u/h}$ twist prefactor}. The
prefactor does not cancel between numerator and denominator: it contributes an overall
$\kappa_1^{s_1+s_2}$ (analytically, $Q_1(\theta_i+sh)/Q_1(\theta_i)=\kappa_1^{s}\,
q_1(\theta_i+sh)/q_1(\theta_i)$ since $h=1$). Using the bare polynomial $q_1$ in place of $Q_1$ in
\eqref{eq:wavefunction} is \emph{not} equivalent and fails numerically by an order-one amount.
\emph{Verification:} literal $0$ residual over $297$ cases (41 states $\times$ every $(s_1,s_2)$
pair; Section K, cells 139--141).

\section{Main result: the closed two-rotation identity}
\label{sec:main}

Substituting \eqref{eq:gamma} and \eqref{eq:A-decomp}/\eqref{eq:B-diag} into the expansion of
\eqref{eq:Odef} in powers of $E_{21}^a\otimes E_{21}^b$ and using \eqref{eq:wavefunction} for every
$x_2$-pairing gives the fully assembled, closed formula:
\begin{equation}
  \boxed{\;
  O(k_1,k_2) \;=\; \sum_{a=0}^{\lambda_1}\sum_{b=0}^{\lambda_2} c_A(a;k_1)\,c_B(b;k_2)
  \!\!\sum_{s+t=a+b}\!\! B(a,b;s,t)\,W(s,t)\;}
  \label{eq:main}
\end{equation}
\begin{equation}
  c_A(a;k_1) := \frac{(-k_1)^a}{a!}\,\gamma_a(\lambda_1), \qquad
  c_B(b;k_2) := \frac{(-k_2)^b}{b!}\,\gamma_b(\lambda_2),
  \label{eq:cAcB}
\end{equation}
with $\gamma$ as in \eqref{eq:gamma}, $B$ as in \S\ref{sec:recoupling}, and $W$ as in
\eqref{eq:wavefunction}. \emph{Verification:} literal $0$ residual across the full 41-state,
six-representation sweep, on two independent $(k_1,k_2)$ grid points per state ($82$ cases total),
directly against the independently-defined operatorial object \eqref{eq:Odef} (Section K, cells
142--151). At $k_2=0$, \eqref{eq:main} reduces exactly to the single-rotation dictionary
\eqref{eq:single-dict}. \eqref{eq:main} is, to our knowledge, the first fully closed statement of
the two-independent-local-rotation overlap for the $L=2$ XXX chain as an explicit finite sum of
Baxter Q-function bilinears.

\section{The derived sum rule}
\label{sec:sumrule}

Because $\mu_i$'s poles are simple,
\[
  \mathrm{br}_i[f]=\sum_s w_i(s)\,f(\theta_i+sh), \qquad
  w_i(s):=2\pi i\,\mathrm{Res}_{u=\theta_i+sh}\mu_i(u),
\]
pointwise (\S\ref{sec:setup}). Expanding
\eqref{eq:det}'s $2\times2$ determinant this way and using \eqref{eq:wavefunction} gives $O(k,k)$ as
an explicit Q-bilinear sum with known coefficients,
\begin{multline}
  \det[\lambda_1,\lambda_2]\bigl[k^{u/h}Q_1(u)\bigr] \;=\; \\
  \sum_{s,t} w_1(s)\,w_2(t)\,\frac{\theta_2-\theta_1+(t-s)h}{\theta_2-\theta_1}\;
  k^{s+t}\;W(s,t),
  \label{eq:det-expanded}
\end{multline}
verified exactly against the live $\det[\cdot]$ object. Equating \eqref{eq:det-expanded} to
\eqref{eq:main} at $k_1=k_2=k$, term-by-term in $k^{s+t}\,W(s,t)$, yields the sum rule on $B$:
\begin{equation}
  \boxed{\;(-1)^{s+t}\!\!\sum_{a+b=s+t}\!\! \frac{\gamma_a(\lambda_1)\,\gamma_b(\lambda_2)}{a!\,b!}\,
  B(a,b;s,t) \;=\; w_1(s)\,w_2(t)\,\frac{\theta_2-\theta_1+(t-s)h}{\theta_2-\theta_1}\;}
  \label{eq:sumrule}
\end{equation}
\emph{Verification:} literal $0$ deviation over $20$ cases across three representations (Section K,
cells 146--148).

\section{Closing remarks}
\label{sec:closing}

\textbf{Structural verdict.} The two-rotation data is carried entirely by the recoupling matrix $B$
(equivalently $A^{-1}$), not by any spectral-parameter insertion into the FSoV bracket. The
pole-cancelling polynomials of \S\ref{sec:pole-cancel} localize a bracket to a single tower exactly,
but they cannot \emph{graded-select} a single level within that tower: no naive per-row-insertion
determinant (twisting each site's bracket independently by its own $k_\alpha^{u/h}$, with node
normalizations fixed to match \eqref{eq:step0} at $k_1=k_2$) reproduces \eqref{eq:main} --- its
$k_1^ak_2^b$ coefficient is always the single "diagonal" term $R(a,b)W(a,b)$ (with
$R(s,t):=w_1(s)w_2(t)\frac{\theta_2-\theta_1+(t-s)h}{\theta_2-\theta_1}$ as in
\eqref{eq:det-expanded}), whereas the true coefficient is the full $B$-weighted sum over the level.
The gap between the two is exactly the physics $A/B$ encode; it is not a small correction on top of
a naive determinant.

\textbf{Open.} The off-diagonal closed form of $A$ (hence $B$) beyond \eqref{eq:A-diag}/\eqref{eq:B-diag}
was not found; the growing complexity with level suggests a discrete-orthogonal-polynomial family
(Hahn or Racah type) rather than a simple product formula. Exact rational values are available in
the notebook for every entry with $\lambda_1,\lambda_2\le3$.

\textbf{Next target.} The natural generalization is the analogue of $B$ for general
$\mathfrak{gl}(3)$ representations, where the recoupling is expected to become genuinely
multi-index and the connection to known 6j/Racah-type combinatorics correspondingly sharper.

The full verification record for this program, including negative results (failed candidate
constructions, an incorrect conjectured sum rule later corrected, and the derivation history behind
each boxed result above), lives in \texttt{XXX\_CG\_L2\_V2.wb}, Sections G--K (cells 78--154).

\end{document}
